\subsection{Current Limitations}\label{subsec:current-limitations}

\textbf{On-pod cryptography:} Current version uses gateway-based \ac{AEAD} encryption.
Pod payloads are encrypted in transit (\ac{LoRa} frames) but encryption keys are managed by gateway.
Future work will implement on-pod key derivation for end-to-end encryption.

\textbf{Key rotation:} \ac{ADR}-001 specifies key rotation workflows, but automation is not yet implemented.
Manual key updates are required for long-term deployments (target: automated weekly ratcheting).

\textbf{Replay window size:} Fixed at 64 frames.
For high-latency scenarios (weeks offline), larger windows or epoch-based schemes may be needed.

\textbf{Single gateway and state sync:} Current architecture assumes a single gateway per site.
Multi-gateway deployments with state synchronization and consistent replay-window enforcement are non-trivial and deferred.

\textbf{OTS upgrade automation:} Proofs currently require an \texttt{ots upgrade} lifecycle step that may be manual in constrained environments.
Future versions will add automation and error handling.

% --- New: explicit safety-oriented framing ---
\subsection{Safety-relevant limitations}\label{subsec:safety-relevant-limitations}
TrackOne improves the integrity and auditability of telemetry but does not, in its current deliverable, provide the guarantees required for autonomous actuation in safety-sensitive deployments.
Concretely:
\begin{itemize}[nosep]
  \item The system does not guarantee bounded-latency delivery of facts under all field conditions; decision logic that depends on real-time telemetry must be gated by human-in-the-loop checks (see R-3 in Requirements).
  \item Hardware failure modes (sensor drift, electrode corrosion, ADC saturation) are outside the scope of TrackOne's core pipeline and must be mitigated by hardware test plans and periodic calibration procedures.
  \item No formal \ac{SIL}-level certification exists for control; TrackOne outputs should be treated as auditable evidence, not a certification of safe-to-act status.
\end{itemize}

% --- New: Documentation & Traceability ---
\subsection{Documentation and traceability}\label{subsec:documentation-and-traceability}
For an end-of-year project intended to be deployed in challenging environments, documentation is not optional.
The following artifacts are ``living'' and must be kept in sync with code and deployments:
\begin{itemize}[nosep]
  \item Requirements (including the irrigation-specific requirements added in this paper) with traced links to tests and \ac{ADR}s.
  \item Architecture diagrams and deployment profiles (per-site) that document pod/gateway/backend interactions and escalation paths.
  \item Test plans and runbooks covering on-pod diagnostics, failure-injection scenarios (link loss, clock drift, storage exhaustion), and recovery procedures.
  \item Known limitations and mitigation checklists attached to each release (what an operator must not do with this release).
\end{itemize}

These artifacts should be versioned alongside software releases and available to field operators and auditors.

% --- New: Operational procedures (concise) ---
\subsection{Operational procedures}\label{subsec:operational-procedures}
Deploying TrackOne in irrigation contexts requires clear operational playbooks.
Minimal recommended procedures:
\begin{enumerate}[nosep]
  \item Site bootstrap: verify provisioning (device table v1.0), perform a 72-hour pilot with local monitoring enabled, and confirm anchor delivery for two consecutive days before enabling any automatic actuation that affects irrigation.
  \item Firmware update: stage updates to a pilot subset of devices for 7 days; monitor telemetry integrity and rollback on anomalous behavior.
  \item Storage exhaustion: when on-device or gateway storage reaches 80\% capacity, system must stop non-critical writes and alert operators; at 95\% capacity it must trigger emergency archival and operator intervention.
  \item Operator override logging: any manual override must create an auditable record that is included in the next day's Merkle root and proof, so that later audits can reconstruct operator interventions.
\end{enumerate}

\subsection{Ethical and regulatory considerations}\label{subsec:ethical-and-regulatory-considerations}

\textbf{Fact confidentiality:} Telemetry payloads remain local (gateway storage).\ Only Merkle roots are public (Bitcoin blockchain via \ac{OTS}). Individual facts cannot be reconstructed from blockchain data.

\textbf{Right to be forgotten (GDPR):} Local fact deletion is possible (gateway operator decision).\ However, Merkle root in Bitcoin blockchain is immutable.\ Heritage sites should assess \ac{GDPR} applicability before deployment.

\textbf{Metadata leakage:} Frame timing and pod IDs are visible to \ac{LoRa} eavesdroppers.\ XChaCha20 encryption hides payload content but not traffic patterns.

\textbf{Open standards and reuse:} Publishing a SensorThings-compatible \ac{API} means heritage authorities and researchers can consume TrackOne telemetry with off-the-shelf clients; \ac{GIS} and hydrology feeds can be joined without bespoke adapters.
The canonical TrackOne facts and \ac{OTS}-anchored Merkle roots underpin that \ac{API}, so the public view is a projection over an immutable ledger rather than a mutable database without provenance.

\subsection{Environmental Impact}\label{subsec:environmental-impact}

\textbf{Battery disposal:} CR2032 lithium cells require proper recycling.\ Multi-year autonomy reduces disposal frequency vs.\ annually-replaced batteries.

\textbf{Bitcoin energy consumption:} \ac{OTS} leverages existing Bitcoin infrastructure (no additional mining).\ Daily batching minimizes per-fact environmental impact (1 anchor amortized across 1000+ facts).

\subsection{Heritage Site Stewardship}\label{subsec:heritage-site-stewardship}

\textbf{Reversibility:} Adhesive pod mounts are non-invasive and removable without surface damage.\ Critical for heritage site preservation ethics.

\textbf{Curator control:} Gateway operator decides what data to collect and publish (Merkle roots).\ Public verification enables accountability but respects site autonomy.

\textbf{Long-term access:} Open-source design ensures data remains accessible even if TrackOne project discontinues.\ Bitcoin verification survives organizational changes.

\subsection{Responsible Disclosure}\label{subsec:responsible-disclosure}

\textbf{Security vulnerabilities:} All cryptographic designs documented in \ac{ADR}s. Independent security review encouraged (MIT license).\ No security-by-obscurity.

\textbf{Test coverage:} ~85\% coverage with 181 tests provides confidence but not proof of correctness.

\textbf{Known attack vectors:} Physical pod tampering, gateway compromise, replay attacks beyond window size.\ Documented in \ac{ADR}-002 threat model.

\subsection{Stakeholder expectations and ethical scope management}\label{subsec:stakeholder-expectations-and-ethical-scope-management}
Field stakeholders occasionally push comfort-zone technologies (e.g., generic image watermarking or steganography) even when they do not match the telemetry problem space.
Heritage pods emit scalar facts; there are no images today.
Implementing watermarking to appease expectations would misrepresent delivered functionality and divert scarce energy budget.
Instead, the project documents alternative tracks (Option B QIM watermarking, Option C PyO3/Rust port) in \ac{ADR}s, clearly labeling them as future work contingent on hardware funding and bench instrumentation.

Equipment discussions are treated transparently: new hardware (oscilloscope, FPGA stimulus%, bio-impedance rigs
) must map to a concrete experiment plan tied to documented milestones.
Until funding is committed, TrackOne prioritizes the achievable Option A pipeline and avoids promising features (``watermark'' keywords) that are not implemented.
This preserves integrity of reporting and maintains trust with community partners who rely on verifiable telemetry rather than speculative features.

\subsection{Internet-Draft Submission}
\label{subsec:internet-draft}

As a direct output of this project, we prepared an \ac{IETF} Internet-Draft
(\texttt{draft-trackone-verifiable-telemetry-ledgers-00}) specifying the
verifiable telemetry ledger profile described in this report.
The draft
targets the Independent Submission stream as an Informational document
and was validated with \texttt{xml2rfc}~v3 and \texttt{idnits}.

\subsubsection{Scope and Positioning}

The draft codifies the pipeline semantics developed in this project ---
canonical fact encoding, deterministic Merkle commitment, day chaining,
and OpenTimestamps anchoring --- into a self-contained specification
suitable for independent implementation.
It explicitly positions itself
as complementary to active \ac{IETF} work:

\begin{itemize}
    \item \textbf{SCITT} (draft-ietf-scitt-architecture): generic
    supply-chain transparency; our profile addresses the
    disconnected, daily-batching case S\ac{CI}TT does not cover.
    \item \textbf{COSE Merkle Tree Proofs} (draft-ietf-cose-merkle-tree-proofs):
    \ac{CBOR}-encoded proof primitives; our implementation includes a
    \ac{CBOR} migration path via \texttt{trackone-core}.
    \item \textbf{RATS}: device attestation; deferred in our profile to
    the provisioning layer.
\end{itemize}

\subsubsection{Implementation Status}

The draft references the published Rust crates (\texttt{trackone-core},
\texttt{trackone-ledger}, and others on \texttt{crates.io}) and the
Python gateway as a reference implementation per \ac{RFC}~7942.
The pipeline
has been validated end-to-end on an STM32L476RG development board with
emulated telemetry links; no production field deployment has been
conducted at the time of writing.

\subsubsection{Disclosure}

Early structural drafts of the Internet-Draft were prepared with AI
writing assistance.
All technical content and design decisions reflect
the author's work as documented throughout this report and the project's
architectural decision records.
This is disclosed in the draft's
Acknowledgments section, consistent with \ac{IETF} guidance on AI-assisted
contributions.


%\subsection{Physarum/Khettara metaphors: stewardship, resilience, and planning}
%Design metaphors are not decoration; they shape operational choices. The \emph{Physarum} colony suggests distributed, self-healing sensing with tiny energy budgets; the \emph{khettara} evokes long-lived, low-maintenance infrastructure. TrackOne reflects both by:
%\begin{itemize}[nosep]
%  \item favoring verifiable evidence over heavy online services (peer attestation + daily anchors rather than always-on cloud);
%  \item encoding continuity tools that match field realities (stationary calendar during outages, backlog clearance, immutable sidecars);
%  \item integrating GIS constraints in deployment (buffer zones, access routes, flood sentinels) so resilience is planned, not improvised.
%\end{itemize}
%These metaphors also discipline trade-offs: we avoid invasive mounts and high-frequency telemetry that jeopardize heritage fabric or battery life, and we accept slower, scheduled anchoring in exchange for multi-year autonomy. Stewardship, in this framing, is the commitment to leave an auditable, reversible trace that future curators can verify without depending on a single vendor or server.
